Verfolge ein wenig Wasser. Nullniveau der Lageenergie: die Düse. Zustand 1: das Wasser an der Düse (nur kinetische Energie). Zustand 2: bei halber Geschwindigkeit (Lageenergie und kinetische Energie). Zustand 3: zuoberst (nur Lageenergie).

\begin{abcliste}
\abc
Sei $H$ die Höhe des Strahls und $v_0$ die Geschwindigkeit an der Düse. Zwischen den Zuständen 1 und 3 wird kinetische Energie zu Lageenergie:
\begin{align}
 \frac{1}{2}\,m\,v_0^2 &= m\,g\,H \\
 v_0 &= \vDF \\
   &= \sqrt{2\cdot\g\cdot\Hs} \\
   &= \vD \approx \result{\vDP{0}{2}}
\end{align}

\abc
Im Zustand 2, auf der Höhe $h$, ist die Geschwindigkeit $v_0/2$. Energieerhaltung zwischen den Zuständen 1 und 2, mit $\frac{1}{2}\,m\,v_0^2 = m\,g\,H$:
\begin{align}
 m\,g\,H &= m\,g\,h + \frac{1}{2}\,m\left(\frac{v_0}{2}\right)^2 = m\,g\,h + \frac{1}{4}\,m\,g\,H
\end{align}
Also
\begin{align}
 h &= \hbF \\
   &= \frac{3}{4}\cdot\Hs \\
   &= \hb \approx \result{\hbP{0}{2}}
\end{align}
Halbe Geschwindigkeit heisst ein Viertel der kinetischen Energie: drei Viertel des Wegs nach oben.
\end{abcliste}
